\paragraph{Nhạc âm - Tạp âm}
\begin{itemize}
\item \textbf{Nhạc âm} là âm có tần số xác định và đồ thị dao động là đường tuần hoàn xác định.
\item \textbf{Tạp âm} là âm không có một tần số xác định và đồ thị dao động là những đường cong phức tạp.
\end{itemize}
\paragraph{Âm cơ bản - Họa âm}
\begin{itemize}
\item Họa âm bậc 1 có tần số $f_1 = f_0$
\item Họa âm bậc 2 có tần số $f_2 = 2f_0$
\item Họa âm bậc 3 có tần số $f_3 = 3f_0$
\item ...
\item Họa âm bậc n có tần số $f_n = nf_0$
\end{itemize}
\paragraph{Nguồn nhạc âm}
\begin{itemize}
\item Giải thích sự tạo thành âm do dây dao động: khi trên dây xuất hiện sóng dừng có những chỗ sợi dây dao động với biên độ cực đại (bụng sóng), đẩy không khí xung quanh nó một cách tuần hoàn và do đó phát ra một sóng âm tương đối mạnh có cùng tần số dao động của dây.
\begin{align*}\ell=k\dfrac{\lambda}{2}=k\dfrac{v}{2f}\Rightarrow \text{(với k = 1; 2; 3;...)} \Rightarrow f=k\dfrac{v}{2\ell}\end{align*} 
Tần số âm cơ bản là \sh{$f_1=\dfrac{v}{2\ell},$} họa âm bậc 1 là \sh{$f_2=2.\dfrac{v}{2\ell}=2f_1$}, họa âm bậc 2 là \sh{$f_3=3.\dfrac{v}{2\ell}=3f_1,... $}. 
\item Giải thích sự tạo thành âm do cột không khí dao động:
Khi sóng âm (sóng dọc) truyền qua không khí trong một ống, chúng phản xạ ngược lại ở mỗi đầu và đi trở lại qua ống (sự phản xạ này vẫn xảy ra ngay cả khi đầu để hở). Khi chiều dài của ống phù hợp với bước sóng của sóng âm $\left(\ell=k\dfrac{\lambda}{2}, \ \ct{hoặc}\ \ell=\left(k+\dfrac{1}{2}\right)\dfrac{\lambda}{2}\right)$ thì trong ống xuất hiện sóng dừng sinh ra âm.
\end{itemize}
\paragraph{Hiệu ứng Doppler}
\notebox{\textbf{Hiệu ứng Doppler} là sự thay đổi tần số sóng (mà máy thu ghi nhận được) khi có sự chuyển động tương đối giữa máy thu và nguồn phát}
Gọi v là vận tốc truyền âm, f là tần số của âm.
\begin{itemize}
\begin{itemize}
\item Nguồn âm đứng yên, máy thu chuyển động với vận tốc $v_{M}$.
\begin{itemize}
\item Máy thu chuyển động lại gần nguồn âm thì thu được âm có tần số: $f'=\dfrac{v+v_M}{v} f_0$
\item Máy thu chuyển động ra xa nguồn âm thì thu được âm có tần số: $f'=\dfrac{v-v_M}{v} f_0$
\end{itemize}
\item Nguồn âm chuyển động với vận tốc $v_{S}$, máy thu đứng yên.
\begin{itemize}
\item Nguồn âm chuyển động lại gần máy thu thì thu được âm có tần số: $f'=\dfrac{v}{v-v_S} f_0$
\item Nguồn âm chuyển động ra xa máy thu thì thu được âm có tần số: $f'=\dfrac{v}{v+v_S}f_0$
\end{itemize}
\item Ta có công thức tổng quát xác định tần số mà máy thu thu được:
$f'=\dfrac{v \pm v_M}{v \pm v_S} f_0$
\end{itemize}
